\documentclass[10pt]{beamer}

\usetheme{metropolis}
\usepackage{appendixnumberbeamer}

\usepackage[spanish]{babel}
\usepackage[utf8]{inputenc}

\usepackage{booktabs}
\usepackage[scale=2]{ccicons}

\usepackage{pgfplots}
\usepgfplotslibrary{dateplot}

\usepackage{xspace}
\usepackage{epsfig}
\usepackage{color}
\usepackage{listings}
\usepackage[]{algorithm2e}
\newcommand{\themename}{\textbf{\textsc{metropolis}}\xspace}
\graphicspath{{./data/}{./resources/}}
\definecolor{mygray}{gray}{0.4}

\title{Taller de programación competitiva}
\subtitle{\textit{An amateur approach}}
\date{20 de febrero de 2018}
\author{
  Ignacio Ballesteros González\\
  {\color{mygray}\texttt{ballesteros@acm.org}\\}
  \\
  Samuel García Haro\\
  {\color{mygray}\texttt{samgh96@gmail.com}\\}
}

\institute{}
\titlegraphic{\hfill\includegraphics[height=1.5cm]{acm_png.png}}

\begin{document}

\maketitle

\begin{frame}{Índice}
  \setbeamertemplate{section in toc}[sections numbered]
  \tableofcontents[hideallsubsections]
\end{frame}

\section{Introduction}

\section{Estructuras de Datos básicas}

\section{Entrada/Salida}
\begin{frame}{Entrada/Salida}
  \begin{itemize}
  \item La entrada nos da pistas de la modelización del problema. \pause
  \item Es importante \textbf{no perder más tiempo del necesario con la E/S}. \pause
  \item Se recomienda llevarla apuntada o en su defecto memorizarla.
  \end{itemize}

\end{frame}


\begin{frame}{Entrada/Salida: Entrada}
  \texttt{4 -> número de entradas \newline \pause
    4 -> tamaño de vector \newline \pause
    1 1 1 2 -> elementos de vector \newline \pause
    2 \newline
    1 1 \newline
    5 \newline
    1 1 2 2 3 \newline
  }
\end{frame}

\begin{frame}{Entrada/Salida: Procesador genérico}
  \begin{algorithm}[H]
    \KwData{num\_entradas, tam\_vector, vector}
    declarar variables\;
    \While{num\_entradas \textgreater 0}{
      leer tam\;
      inicializar vector\;
      \For{$i$ in \KwTo tam\_vector}{
        vector[i] = leer\_entero\;
      }
      procesar\;
      num\_entradas -= 1\;
    }
  \end{algorithm}
\end{frame}
\defverbatim[colored]\lstES{
  \begin{lstlisting}[language=Java,basicstyle=\ttfamily,keywordstyle=\color{blue}]
    public static void main(String[] args){
      int n, len;
      int [] arr;
      FastReader fr = new FastReader();

      n = fr.nextInt();
      while (n > 0){
        len = fr.nextInt();
        arr = new int[len];
        for (int i = 0; i < len; i++)
        arr[i] = fr.nextInt();

        System.out.println(findMajority(arr));
        n--;
      }
    }
  \end{lstlisting}
}

\begin{frame}{Entrada/Salida: Ejemplo}
  \lstES
\end{frame}
\section{Métodos algorítmicos}

\subsection{Divide y vencerás}
\begin{frame}{Divide y vencerás: Concepto}
  \begin{itemize}
  \item Patrón de diseño de algoritmos de carácter recursivo. \pause
  \item Descompone problemas en subproblemas de solución más sencilla. \pause
  \item Se distinguen cuatro pasos: \pause
    \begin{itemize}
    \item Hallar el caso base del problema. \pause
    \item Dividir la estructura de datos hasta encontrar el caso base. \pause
    \item Aplicar la solución al caso base. \pause
    \item Recomponer la estructura de datos ya resuelta. \pause
    \end{itemize}
  \item Ejemplos de este método: Mergesort, quicksort, búsqueda binaria...
  \end{itemize}
\end{frame}
\subsection{Método voraz}
\begin{frame}{Método voraz: Concepto}
  \begin{itemize}
  \item Este método algorítmico se basa en la elección de soluciones locales óptimas
    con el fin de obtener soluciones globales óptimas. \pause
  \item Útil en problemas de optimización pero difícil de probar su correctitud. \pause
  \item Consiste en los siguientes fundamentos: \pause
    \begin{itemize}
    \item Disponemos de un conjunto de entrada (candidatos) y de otro de salida. \pause
    \item Mediante heurísticas decidimos qué elemento entra en el conjunto de salida 
      (y lo eliminamos de los candidatos). \pause
    \item Para decidir los siguientes elementos los sometemos a un test de factibilidad. \pause
    \end{itemize}
    \item Ejemplos: A*, Prim, Kruskal, TSP...
  \end{itemize}
\end{frame}
\subsection{Programación dinámica}

\end{document}